% Standalone editable visual-model calculation note for issue 15853.
\section{Period Hat Silhouettes}
All coordinates are unit head lengths in the existing head frame, $x$ forward
and $z$ up. The straw brim retains its earlier closed mesh. The crown and band start from
their original closed shells; for $z>0.82$, put
$t=(z-0.82)/(z_{\max}-0.82)$ and map its height to
$0.82+(z_{\max}-0.82)[0.15t+0.85(1-(1-t)^3)]$. The positive derivative
preserves shell orientation, and both shells share this map to stay attached.
The unchanged lower edge preserves the
brim and band junction. This creates a broad flat top with steeper sides; the
hair clearance is checked by sampled rays. It is a stylized approximation of
period headwear rather than a measured reproduction.

The flat cap and fuller newsboy cap use the same short peak and cropped hair.
The newsboy crown uses offset $0.085$ and forward overhang $0.16$ in head
lengths. The front crown is smoothly pulled toward a central peak point, with the
pull bounded for longer peaks; a small ellipsoid at that join represents the
snap fastening. The band is cut to a lower upper edge before the crown map
so its overall height stays near the original. The baseball-cap path is unchanged. All
parts are visual meshes: none contributes to mass, contact or inertia.
Acceptance requires finite, outward closed meshes, connected snap endpoints,
crown hair clearance and native dynamics identity at representative states.
\begin{verbatim}
PYTHONPATH=.:vendor/ud-tools/src MUJOCO_GL=egl \
  python -m pytest tests/unit/model_appearance/test_period_hats.py
PYTHONPATH=.:vendor/ud-tools/src MUJOCO_GL=egl \
  python scripts/render_period_hats.py --out /tmp/period-hats
\end{verbatim}
